
\documentclass{article}
\usepackage{colortbl}
\usepackage{makecell}
\usepackage{multirow}
\usepackage{supertabular}

\begin{document}

\newcounter{utterance}

\centering \large Interaction Transcript for game `dond', experiment `semi\_en', episode 16 with Qwen3{-}32B{-}t0.0.
\vspace{24pt}

{ \footnotesize  \setcounter{utterance}{1}
\setlength{\tabcolsep}{0pt}
\begin{supertabular}{c@{$\;$}|p{.15\linewidth}@{}p{.15\linewidth}p{.15\linewidth}p{.15\linewidth}p{.15\linewidth}p{.15\linewidth}}
    \# & $\;$A & \multicolumn{4}{c}{Game Master} & $\;\:$B\\
    \hline

    \theutterance \stepcounter{utterance}  
    & & \multicolumn{4}{p{0.6\linewidth}}{
        \cellcolor[rgb]{0.9,0.9,0.9}{
            \makecell[{{p{\linewidth}}}]{
                \texttt{\tiny{[P1$\langle$GM]}}
                \texttt{You are playing a negotiation game in which you have to agree on how to divide a set of items among you and another player.} \\
\\ 
\texttt{Rules:} \\
\texttt{(a) You and the other player are given a set of items. Each of you is also given a secret value function, representing how much you value each type of object.} \\
\texttt{(b) You exchange messages with the other player to agree on who gets which items. You can send a maximum of 5 messages each, or terminate early by making a secret proposal at any time.} \\
\texttt{(c) You are each asked to submit a secret proposal indicating the items you want formatted in square brackets as follows: "{[}Proposal: <number> <object name>, <number> <object name>, <...>{]}"} \\
\texttt{(d) If your proposals are complementary, i.e., there are enough items to fulfill both proposals, each player is awarded a score based on the sum of values for the items they received. Otherwise, both of you get zero points.} \\
\texttt{(e) Your goal is to maximize the score you receive.} \\
\\ 
\texttt{Let us start.} \\
\\ 
\texttt{The set of available items is:} \\
\\ 
\texttt{3 vases, 1 chair, 1 wire} \\
\\ 
\texttt{Your secret value function is:} \\
\\ 
\texttt{vase: 2, chair: 3, wire: 1} \\
\\ 
\texttt{IMPORTANT: Your messages, unless it is the secret proposal, are directly transmitted to the other player, so do not include any response to the rules or text announcing your message. To make a secret proposal, use the indicated format. Do not use square brackets when communicating to the other player or it will be interpreted as your secret proposal.} \\
            }
        }
    }
    & \\ \\

    \theutterance \stepcounter{utterance}  
    & \multicolumn{4}{p{0.6\linewidth}}{
        \cellcolor[rgb]{0.9,0.9,0.9}{
            \makecell[{{p{\linewidth}}}]{
                \texttt{\tiny{[P1$\rangle$GM]}}
                \texttt{<think>} \\
\texttt{Okay, let's see. I need to negotiate the division of 3 vases, 1 chair, and 1 wire with another player. My values are vase: 2, chair: 3, wire: 1. The other player's values are secret, so I don't know what they value. My goal is to maximize my own score.} \\
\\ 
\texttt{First, I should consider what the other player might want. Since there are only 3 vases, maybe they value vases highly. But the chair and wire are unique, so if the other player values either of those, they might want them. But since I don't know their values, I need to make some assumptions or try to get information.} \\
\\ 
\texttt{But since I can't ask them directly, maybe I should propose a division that gives me the chair and some vases, assuming that the other player might prefer the wire and the remaining vases. Alternatively, if they want the chair, I might have to give it up. But since my chair is worth 3, which is higher than the vases (2 each) and wire (1), I should prioritize getting the chair if possible.} \\
\\ 
\texttt{Wait, but there are 3 vases. If I take the chair, maybe I can take 2 vases, and leave 1 vase and the wire for the other player. Let me calculate my score in that case: chair (3) + 2 vases (2*2=4) = 7. The other player would get 1 vase (2) + wire (1) = 3. But if their values are such that they value the wire more, maybe that's acceptable. But if they value vases more, maybe they would prefer more vases. But since I don't know, I need to make a proposal that is likely to be complementary.} \\
\\ 
\texttt{Alternatively, maybe I should take the chair and 1 vase, leaving 2 vases and the wire. My score would be 3 + 2 = 5. The other player could get 2 vases (4) + wire (1) = 5. That's a total of 10, which is better than the previous 7+3=10. Wait, same total. But maybe the other player would prefer that split. But since I don't know their values, maybe splitting the vases more evenly is better. But I value the chair more, so} \\
            }
        }
    }
    & & \\ \\

    \theutterance \stepcounter{utterance}  
    & & & \multicolumn{2}{p{0.3\linewidth}}{
        \cellcolor[rgb]{0.9,0.9,0.9}{
            \makecell[{{p{\linewidth}}}]{
                \texttt{\tiny{[GM$|$GM]}}
                \texttt{continue} \\
            }
        }
    }
    & & \\ \\

    \theutterance \stepcounter{utterance}  
    & & & \multicolumn{2}{p{0.3\linewidth}}{
        \cellcolor[rgb]{0.9,0.9,0.9}{
            \makecell[{{p{\linewidth}}}]{
                \texttt{\tiny{[GM$|$GM]}}
                \texttt{<think>} \\
\texttt{Okay, let's see. I need to negotiate the division of 3 vases, 1 chair, and 1 wire with another player. My values are vase: 2, chair: 3, wire: 1. The other player's values are secret, so I don't know what they value. My goal is to maximize my own score.} \\
\\ 
\texttt{First, I should consider what the other player might want. Since there are only 3 vases, maybe they value vases highly. But the chair and wire are unique, so if the other player values either of those, they might want them. But since I don't know their values, I need to make some assumptions or try to get information.} \\
\\ 
\texttt{But since I can't ask them directly, maybe I should propose a division that gives me the chair and some vases, assuming that the other player might prefer the wire and the remaining vases. Alternatively, if they want the chair, I might have to give it up. But since my chair is worth 3, which is higher than the vases (2 each) and wire (1), I should prioritize getting the chair if possible.} \\
\\ 
\texttt{Wait, but there are 3 vases. If I take the chair, maybe I can take 2 vases, and leave 1 vase and the wire for the other player. Let me calculate my score in that case: chair (3) + 2 vases (2*2=4) = 7. The other player would get 1 vase (2) + wire (1) = 3. But if their values are such that they value the wire more, maybe that's acceptable. But if they value vases more, maybe they would prefer more vases. But since I don't know, I need to make a proposal that is likely to be complementary.} \\
\\ 
\texttt{Alternatively, maybe I should take the chair and 1 vase, leaving 2 vases and the wire. My score would be 3 + 2 = 5. The other player could get 2 vases (4) + wire (1) = 5. That's a total of 10, which is better than the previous 7+3=10. Wait, same total. But maybe the other player would prefer that split. But since I don't know their values, maybe splitting the vases more evenly is better. But I value the chair more, so} \\
            }
        }
    }
    & & \\ \\

    \theutterance \stepcounter{utterance}  
    & & \multicolumn{4}{p{0.6\linewidth}}{
        \cellcolor[rgb]{0.9,0.9,0.9}{
            \makecell[{{p{\linewidth}}}]{
                \texttt{\tiny{[GM$\rangle$P2]}}
                \texttt{You are playing a negotiation game in which you have to agree on how to divide a set of items among you and another player.} \\
\\ 
\texttt{Rules:} \\
\texttt{(a) You and the other player are given a set of items. Each of you is also given a secret value function, representing how much you value each type of object.} \\
\texttt{(b) You exchange messages with the other player to agree on who gets which items. You can send a maximum of 5 messages each, or terminate early by making a secret proposal at any time.} \\
\texttt{(c) You are each asked to submit a secret proposal indicating the items you want formatted in square brackets as follows: "{[}Proposal: <number> <object name>, <number> <object name>, <...>{]}"} \\
\texttt{(d) If your proposals are complementary, i.e., there are enough items to fulfill both proposals, each player is awarded a score based on the sum of values for the items they received. Otherwise, both of you get zero points.} \\
\texttt{(e) Your goal is to maximize the score you receive.} \\
\\ 
\texttt{Let us start.} \\
\\ 
\texttt{The set of available items is:} \\
\\ 
\texttt{3 vases, 1 chair, 1 wire} \\
\\ 
\texttt{Your secret value function is:} \\
\\ 
\texttt{vase: 0, chair: 8, wire: 2} \\
\\ 
\texttt{IMPORTANT: Your messages, unless it is the secret proposal, are directly transmitted to the other player, so do not include any response to the rules or text announcing your message. To make a secret proposal, use the indicated format. Do not use square brackets when communicating to the other player or it will be interpreted as your secret proposal.} \\
\\ 
\\ 
\texttt{<think>} \\
\texttt{Okay, let's see. I need to negotiate the division of 3 vases, 1 chair, and 1 wire with another player. My values are vase: 2, chair: 3, wire: 1. The other player's values are secret, so I don't know what they value. My goal is to maximize my own score.} \\
\\ 
\texttt{First, I should consider what the other player might want. Since there are only 3 vases, maybe they value vases highly. But the chair and wire are unique, so if the other player values either of those, they might want them. But since I don't know their values, I need to make some assumptions or try to get information.} \\
\\ 
\texttt{But since I can't ask them directly, maybe I should propose a division that gives me the chair and some vases, assuming that the other player might prefer the wire and the remaining vases. Alternatively, if they want the chair, I might have to give it up. But since my chair is worth 3, which is higher than the vases (2 each) and wire (1), I should prioritize getting the chair if possible.} \\
\\ 
\texttt{Wait, but there are 3 vases. If I take the chair, maybe I can take 2 vases, and leave 1 vase and the wire for the other player. Let me calculate my score in that case: chair (3) + 2 vases (2*2=4) = 7. The other player would get 1 vase (2) + wire (1) = 3. But if their values are such that they value the wire more, maybe that's acceptable. But if they value vases more, maybe they would prefer more vases. But since I don't know, I need to make a proposal that is likely to be complementary.} \\
\\ 
\texttt{Alternatively, maybe I should take the chair and 1 vase, leaving 2 vases and the wire. My score would be 3 + 2 = 5. The other player could get 2 vases (4) + wire (1) = 5. That's a total of 10, which is better than the previous 7+3=10. Wait, same total. But maybe the other player would prefer that split. But since I don't know their values, maybe splitting the vases more evenly is better. But I value the chair more, so} \\
            }
        }
    }
    & \\ \\

    \theutterance \stepcounter{utterance}  
    & & & \multicolumn{4}{p{0.6\linewidth}}{
        \cellcolor[rgb]{0.9,0.9,0.9}{
            \makecell[{{p{\linewidth}}}]{
                \texttt{\tiny{[GM$\langle$P2]}}
                \texttt{<think>} \\
\texttt{Okay, let's try to figure this out. The user is playing a negotiation game where they need to divide items between themselves and another player. The items are 3 vases, 1 chair, and 1 wire. The user's value function is vase: 2, chair: 3, wire: 1. The other player's values are secret. The goal is to maximize their own score by making a secret proposal that's complementary to the other player's.} \\
\\ 
\texttt{First, the user needs to consider possible strategies. Since the other player's values are unknown, the user has to make assumptions or find a proposal that's likely to be complementary regardless of the other's preferences. The key is to ensure that the proposals from both players don't conflict, so they both get points.} \\
\\ 
\texttt{The user's highest value item is the chair (3), followed by vases (2 each), and the wire (1). The chair is unique, so if the user takes it, the other player can't. The wire is also unique. The vases are three, so there's more flexibility there.} \\
\\ 
\texttt{If the user proposes taking the chair and some vases, they need to leave enough vases and the wire for the other player. Let's see:} \\
\\ 
\texttt{Option 1: User takes chair (3) + 2 vases (2*2=4). Total value: 7. Leaves 1 vase and wire. The other player would need to take those. If their values for vase and wire are such that they can get a positive score, that's good. But if the other player values the chair or vases more, maybe they would prefer that.} \\
\\ 
\texttt{Option 2: User takes chair (3) + 1 vase (2). Total: 5. Leaves 2 vases and wire. The other player could take 2 vases and wire. If the other player values vases higher, this might be better for them. The user's score is lower here, but maybe the other player's proposal would be more complementary.} \\
\\ 
\texttt{Alternatively, if the user takes the wire and some vases, but since the wire is only worth 1 to them, that's probably not optimal. The chair is more valuable.} \\
\\ 
\texttt{Another angle: The other player might also want the chair. If the user takes the chair, the other player can't. But if the other player's value for the chair is higher than the user's, maybe they would} \\
            }
        }
    }
     \\ \\

    \theutterance \stepcounter{utterance}  
    & & & \multicolumn{2}{p{0.3\linewidth}}{
        \cellcolor[rgb]{0.9,0.9,0.9}{
            \makecell[{{p{\linewidth}}}]{
                \texttt{\tiny{[GM$|$GM]}}
                \texttt{continue} \\
            }
        }
    }
    & & \\ \\

    \theutterance \stepcounter{utterance}  
    & & & \multicolumn{2}{p{0.3\linewidth}}{
        \cellcolor[rgb]{0.9,0.9,0.9}{
            \makecell[{{p{\linewidth}}}]{
                \texttt{\tiny{[GM$|$GM]}}
                \texttt{<think>} \\
\texttt{Okay, let's try to figure this out. The user is playing a negotiation game where they need to divide items between themselves and another player. The items are 3 vases, 1 chair, and 1 wire. The user's value function is vase: 2, chair: 3, wire: 1. The other player's values are secret. The goal is to maximize their own score by making a secret proposal that's complementary to the other player's.} \\
\\ 
\texttt{First, the user needs to consider possible strategies. Since the other player's values are unknown, the user has to make assumptions or find a proposal that's likely to be complementary regardless of the other's preferences. The key is to ensure that the proposals from both players don't conflict, so they both get points.} \\
\\ 
\texttt{The user's highest value item is the chair (3), followed by vases (2 each), and the wire (1). The chair is unique, so if the user takes it, the other player can't. The wire is also unique. The vases are three, so there's more flexibility there.} \\
\\ 
\texttt{If the user proposes taking the chair and some vases, they need to leave enough vases and the wire for the other player. Let's see:} \\
\\ 
\texttt{Option 1: User takes chair (3) + 2 vases (2*2=4). Total value: 7. Leaves 1 vase and wire. The other player would need to take those. If their values for vase and wire are such that they can get a positive score, that's good. But if the other player values the chair or vases more, maybe they would prefer that.} \\
\\ 
\texttt{Option 2: User takes chair (3) + 1 vase (2). Total: 5. Leaves 2 vases and wire. The other player could take 2 vases and wire. If the other player values vases higher, this might be better for them. The user's score is lower here, but maybe the other player's proposal would be more complementary.} \\
\\ 
\texttt{Alternatively, if the user takes the wire and some vases, but since the wire is only worth 1 to them, that's probably not optimal. The chair is more valuable.} \\
\\ 
\texttt{Another angle: The other player might also want the chair. If the user takes the chair, the other player can't. But if the other player's value for the chair is higher than the user's, maybe they would} \\
            }
        }
    }
    & & \\ \\

    \theutterance \stepcounter{utterance}  
    & & \multicolumn{4}{p{0.6\linewidth}}{
        \cellcolor[rgb]{0.9,0.9,0.9}{
            \makecell[{{p{\linewidth}}}]{
                \texttt{\tiny{[P1$\langle$GM]}}
                \texttt{<think>} \\
\texttt{Okay, let's try to figure this out. The user is playing a negotiation game where they need to divide items between themselves and another player. The items are 3 vases, 1 chair, and 1 wire. The user's value function is vase: 2, chair: 3, wire: 1. The other player's values are secret. The goal is to maximize their own score by making a secret proposal that's complementary to the other player's.} \\
\\ 
\texttt{First, the user needs to consider possible strategies. Since the other player's values are unknown, the user has to make assumptions or find a proposal that's likely to be complementary regardless of the other's preferences. The key is to ensure that the proposals from both players don't conflict, so they both get points.} \\
\\ 
\texttt{The user's highest value item is the chair (3), followed by vases (2 each), and the wire (1). The chair is unique, so if the user takes it, the other player can't. The wire is also unique. The vases are three, so there's more flexibility there.} \\
\\ 
\texttt{If the user proposes taking the chair and some vases, they need to leave enough vases and the wire for the other player. Let's see:} \\
\\ 
\texttt{Option 1: User takes chair (3) + 2 vases (2*2=4). Total value: 7. Leaves 1 vase and wire. The other player would need to take those. If their values for vase and wire are such that they can get a positive score, that's good. But if the other player values the chair or vases more, maybe they would prefer that.} \\
\\ 
\texttt{Option 2: User takes chair (3) + 1 vase (2). Total: 5. Leaves 2 vases and wire. The other player could take 2 vases and wire. If the other player values vases higher, this might be better for them. The user's score is lower here, but maybe the other player's proposal would be more complementary.} \\
\\ 
\texttt{Alternatively, if the user takes the wire and some vases, but since the wire is only worth 1 to them, that's probably not optimal. The chair is more valuable.} \\
\\ 
\texttt{Another angle: The other player might also want the chair. If the user takes the chair, the other player can't. But if the other player's value for the chair is higher than the user's, maybe they would} \\
            }
        }
    }
    & \\ \\

    \theutterance \stepcounter{utterance}  
    & \multicolumn{4}{p{0.6\linewidth}}{
        \cellcolor[rgb]{0.9,0.9,0.9}{
            \makecell[{{p{\linewidth}}}]{
                \texttt{\tiny{[P1$\rangle$GM]}}
                \texttt{<think>} \\
\texttt{Okay, let's try to figure out the best way to approach this. The user has to divide 3 vases, 1 chair, and 1 wire. Their values are vase: 2, chair: 3, wire: 1. The other player's values are secret, so I need to make a proposal that's likely to be complementary regardless of what the other player wants.} \\
\\ 
\texttt{First, the chair is the highest value item for the user (3 points). If I can get the chair, that's a big boost. But the chair is unique, so if I take it, the other player can't. The wire is only worth 1 to me, so maybe I should leave that for them if they value it more. The vases are three, so there's more flexibility there.} \\
\\ 
\texttt{If I take the chair and two vases, that's 3 + 2*2 = 7 points. Then the other player would have one vase and the wire. If their values for vase and wire are such that they can get a positive score, that's good. But if they value the chair more, maybe they'd want it, but since I took it, their proposal might conflict. However, since I don't know their values, I have to assume they might want the wire or some vases.} \\
\\ 
\texttt{Alternatively, if I take the chair and one vase (3 + 2 = 5), leaving two vases and the wire. The other player could take two vases and the wire. If they value vases more, that's better for them. My score is lower here, but maybe the other player's proposal would be more complementary. But since my goal is to maximize my own score, taking more vases might be better.} \\
\\ 
\texttt{Wait, but the key is that the proposals need to be complementary. If I take the chair and two vases, the other player can only get one vase and the wire. If their proposal requires the chair, which I have, then our proposals would conflict. But since I don't know their values, maybe it's safer to take the chair and one vase, leaving two vases and the wire. That way, even if the other player wants the chair, they can't get it, but maybe they prefer the wire and vases. Alternatively, if they value the wire highly, taking the wire might be better for them.} \\
\\ 
\texttt{Another angle: The wire is only one, so} \\
            }
        }
    }
    & & \\ \\

    \theutterance \stepcounter{utterance}  
    & & & \multicolumn{2}{p{0.3\linewidth}}{
        \cellcolor[rgb]{0.9,0.9,0.9}{
            \makecell[{{p{\linewidth}}}]{
                \texttt{\tiny{[GM$|$GM]}}
                \texttt{continue} \\
            }
        }
    }
    & & \\ \\

    \theutterance \stepcounter{utterance}  
    & & & \multicolumn{2}{p{0.3\linewidth}}{
        \cellcolor[rgb]{0.9,0.9,0.9}{
            \makecell[{{p{\linewidth}}}]{
                \texttt{\tiny{[GM$|$GM]}}
                \texttt{<think>} \\
\texttt{Okay, let's try to figure out the best way to approach this. The user has to divide 3 vases, 1 chair, and 1 wire. Their values are vase: 2, chair: 3, wire: 1. The other player's values are secret, so I need to make a proposal that's likely to be complementary regardless of what the other player wants.} \\
\\ 
\texttt{First, the chair is the highest value item for the user (3 points). If I can get the chair, that's a big boost. But the chair is unique, so if I take it, the other player can't. The wire is only worth 1 to me, so maybe I should leave that for them if they value it more. The vases are three, so there's more flexibility there.} \\
\\ 
\texttt{If I take the chair and two vases, that's 3 + 2*2 = 7 points. Then the other player would have one vase and the wire. If their values for vase and wire are such that they can get a positive score, that's good. But if they value the chair more, maybe they'd want it, but since I took it, their proposal might conflict. However, since I don't know their values, I have to assume they might want the wire or some vases.} \\
\\ 
\texttt{Alternatively, if I take the chair and one vase (3 + 2 = 5), leaving two vases and the wire. The other player could take two vases and the wire. If they value vases more, that's better for them. My score is lower here, but maybe the other player's proposal would be more complementary. But since my goal is to maximize my own score, taking more vases might be better.} \\
\\ 
\texttt{Wait, but the key is that the proposals need to be complementary. If I take the chair and two vases, the other player can only get one vase and the wire. If their proposal requires the chair, which I have, then our proposals would conflict. But since I don't know their values, maybe it's safer to take the chair and one vase, leaving two vases and the wire. That way, even if the other player wants the chair, they can't get it, but maybe they prefer the wire and vases. Alternatively, if they value the wire highly, taking the wire might be better for them.} \\
\\ 
\texttt{Another angle: The wire is only one, so} \\
            }
        }
    }
    & & \\ \\

    \theutterance \stepcounter{utterance}  
    & & \multicolumn{4}{p{0.6\linewidth}}{
        \cellcolor[rgb]{0.9,0.9,0.9}{
            \makecell[{{p{\linewidth}}}]{
                \texttt{\tiny{[GM$\rangle$P2]}}
                \texttt{<think>} \\
\texttt{Okay, let's try to figure out the best way to approach this. The user has to divide 3 vases, 1 chair, and 1 wire. Their values are vase: 2, chair: 3, wire: 1. The other player's values are secret, so I need to make a proposal that's likely to be complementary regardless of what the other player wants.} \\
\\ 
\texttt{First, the chair is the highest value item for the user (3 points). If I can get the chair, that's a big boost. But the chair is unique, so if I take it, the other player can't. The wire is only worth 1 to me, so maybe I should leave that for them if they value it more. The vases are three, so there's more flexibility there.} \\
\\ 
\texttt{If I take the chair and two vases, that's 3 + 2*2 = 7 points. Then the other player would have one vase and the wire. If their values for vase and wire are such that they can get a positive score, that's good. But if they value the chair more, maybe they'd want it, but since I took it, their proposal might conflict. However, since I don't know their values, I have to assume they might want the wire or some vases.} \\
\\ 
\texttt{Alternatively, if I take the chair and one vase (3 + 2 = 5), leaving two vases and the wire. The other player could take two vases and the wire. If they value vases more, that's better for them. My score is lower here, but maybe the other player's proposal would be more complementary. But since my goal is to maximize my own score, taking more vases might be better.} \\
\\ 
\texttt{Wait, but the key is that the proposals need to be complementary. If I take the chair and two vases, the other player can only get one vase and the wire. If their proposal requires the chair, which I have, then our proposals would conflict. But since I don't know their values, maybe it's safer to take the chair and one vase, leaving two vases and the wire. That way, even if the other player wants the chair, they can't get it, but maybe they prefer the wire and vases. Alternatively, if they value the wire highly, taking the wire might be better for them.} \\
\\ 
\texttt{Another angle: The wire is only one, so} \\
            }
        }
    }
    & \\ \\

    \theutterance \stepcounter{utterance}  
    & & & \multicolumn{4}{p{0.6\linewidth}}{
        \cellcolor[rgb]{0.9,0.9,0.9}{
            \makecell[{{p{\linewidth}}}]{
                \texttt{\tiny{[GM$\langle$P2]}}
                \texttt{<think>} \\
\texttt{Okay, let's try to figure out the best way to approach this. The user has to divide 3 vases, 1 chair, and 1 wire. Their values are vase: 2, chair: 3, wire: 1. The other player's values are secret, so I need to make a proposal that's likely to be complementary regardless of what the other player wants.} \\
\\ 
\texttt{First, the chair is the highest value item for the user (3 points). If I can get the chair, that's a big boost. But the chair is unique, so if I take it, the other player can't. The wire is only worth 1 to me, so maybe I should leave that for them if they value it more. The vases are three, so there's more flexibility there.} \\
\\ 
\texttt{If I take the chair and two vases, that's 3 + 2*2 = 7 points. Then the other player would have one vase and the wire. If their values for vase and wire are such that they can get a positive score, that's good. But if they value the chair more, maybe they'd want it, but since I took it, their proposal might conflict. However, since I don't know their values, I have to assume they might want the wire or some vases.} \\
\\ 
\texttt{Alternatively, if I take the chair and one vase (3 + 2 = 5), leaving two vases and the wire. The other player could take two vases and the wire. If they value vases more, that's better for them. My score is lower here, but maybe the other player's proposal would be more complementary. But since my goal is to maximize my own score, taking more vases might be better.} \\
\\ 
\texttt{Wait, but the key is that the proposals need to be complementary. If I take the chair and two vases, the other player can only get one vase and the wire. If their proposal requires the chair, which I have, then our proposals would conflict. But since I don't know their values, maybe it's safer to take the chair and one vase, leaving two vases and the wire. That way, even if the other player wants the chair, they can't get it, but maybe they prefer the wire and vases. Alternatively, if they value the wire highly, taking the wire might be better for them.} \\
\\ 
\texttt{Another angle: The wire is only one, so} \\
            }
        }
    }
     \\ \\

    \theutterance \stepcounter{utterance}  
    & & & \multicolumn{2}{p{0.3\linewidth}}{
        \cellcolor[rgb]{0.9,0.9,0.9}{
            \makecell[{{p{\linewidth}}}]{
                \texttt{\tiny{[GM$|$GM]}}
                \texttt{continue} \\
            }
        }
    }
    & & \\ \\

    \theutterance \stepcounter{utterance}  
    & & & \multicolumn{2}{p{0.3\linewidth}}{
        \cellcolor[rgb]{0.9,0.9,0.9}{
            \makecell[{{p{\linewidth}}}]{
                \texttt{\tiny{[GM$|$GM]}}
                \texttt{<think>} \\
\texttt{Okay, let's try to figure out the best way to approach this. The user has to divide 3 vases, 1 chair, and 1 wire. Their values are vase: 2, chair: 3, wire: 1. The other player's values are secret, so I need to make a proposal that's likely to be complementary regardless of what the other player wants.} \\
\\ 
\texttt{First, the chair is the highest value item for the user (3 points). If I can get the chair, that's a big boost. But the chair is unique, so if I take it, the other player can't. The wire is only worth 1 to me, so maybe I should leave that for them if they value it more. The vases are three, so there's more flexibility there.} \\
\\ 
\texttt{If I take the chair and two vases, that's 3 + 2*2 = 7 points. Then the other player would have one vase and the wire. If their values for vase and wire are such that they can get a positive score, that's good. But if they value the chair more, maybe they'd want it, but since I took it, their proposal might conflict. However, since I don't know their values, I have to assume they might want the wire or some vases.} \\
\\ 
\texttt{Alternatively, if I take the chair and one vase (3 + 2 = 5), leaving two vases and the wire. The other player could take two vases and the wire. If they value vases more, that's better for them. My score is lower here, but maybe the other player's proposal would be more complementary. But since my goal is to maximize my own score, taking more vases might be better.} \\
\\ 
\texttt{Wait, but the key is that the proposals need to be complementary. If I take the chair and two vases, the other player can only get one vase and the wire. If their proposal requires the chair, which I have, then our proposals would conflict. But since I don't know their values, maybe it's safer to take the chair and one vase, leaving two vases and the wire. That way, even if the other player wants the chair, they can't get it, but maybe they prefer the wire and vases. Alternatively, if they value the wire highly, taking the wire might be better for them.} \\
\\ 
\texttt{Another angle: The wire is only one, so} \\
            }
        }
    }
    & & \\ \\

    \theutterance \stepcounter{utterance}  
    & & \multicolumn{4}{p{0.6\linewidth}}{
        \cellcolor[rgb]{0.9,0.9,0.9}{
            \makecell[{{p{\linewidth}}}]{
                \texttt{\tiny{[P1$\langle$GM]}}
                \texttt{<think>} \\
\texttt{Okay, let's try to figure out the best way to approach this. The user has to divide 3 vases, 1 chair, and 1 wire. Their values are vase: 2, chair: 3, wire: 1. The other player's values are secret, so I need to make a proposal that's likely to be complementary regardless of what the other player wants.} \\
\\ 
\texttt{First, the chair is the highest value item for the user (3 points). If I can get the chair, that's a big boost. But the chair is unique, so if I take it, the other player can't. The wire is only worth 1 to me, so maybe I should leave that for them if they value it more. The vases are three, so there's more flexibility there.} \\
\\ 
\texttt{If I take the chair and two vases, that's 3 + 2*2 = 7 points. Then the other player would have one vase and the wire. If their values for vase and wire are such that they can get a positive score, that's good. But if they value the chair more, maybe they'd want it, but since I took it, their proposal might conflict. However, since I don't know their values, I have to assume they might want the wire or some vases.} \\
\\ 
\texttt{Alternatively, if I take the chair and one vase (3 + 2 = 5), leaving two vases and the wire. The other player could take two vases and the wire. If they value vases more, that's better for them. My score is lower here, but maybe the other player's proposal would be more complementary. But since my goal is to maximize my own score, taking more vases might be better.} \\
\\ 
\texttt{Wait, but the key is that the proposals need to be complementary. If I take the chair and two vases, the other player can only get one vase and the wire. If their proposal requires the chair, which I have, then our proposals would conflict. But since I don't know their values, maybe it's safer to take the chair and one vase, leaving two vases and the wire. That way, even if the other player wants the chair, they can't get it, but maybe they prefer the wire and vases. Alternatively, if they value the wire highly, taking the wire might be better for them.} \\
\\ 
\texttt{Another angle: The wire is only one, so} \\
            }
        }
    }
    & \\ \\

    \theutterance \stepcounter{utterance}  
    & \multicolumn{4}{p{0.6\linewidth}}{
        \cellcolor[rgb]{0.9,0.9,0.9}{
            \makecell[{{p{\linewidth}}}]{
                \texttt{\tiny{[P1$\rangle$GM]}}
                \texttt{<think>} \\
\texttt{Okay, let's try to figure out the best way to approach this. The user has to divide 3 vases, 1 chair, and 1 wire. Their values are vase: 2, chair: 3, wire: 1. The other player's values are secret, so I need to make a proposal that's likely to be complementary regardless of what the other player wants.} \\
\\ 
\texttt{First, the chair is the highest value item for the user (3 points). If I can get the chair, that's a big boost. But the chair is unique, so if I take it, the other player can't. The wire is only worth 1 to me, so maybe I should leave that for them if they value it more. The vases are three, so there's more flexibility there.} \\
\\ 
\texttt{If I take the chair and two vases, that's 3 + 2*2 = 7 points. Then the other player would have one vase and the wire. If their values for vase and wire are such that they can get a positive score, that's good. But if they value the chair more, maybe they'd want it, but since I took it, their proposal might conflict. However, since I don't know their values, I have to assume they might want the wire or some vases.} \\
\\ 
\texttt{Alternatively, if I take the chair and one vase (3 + 2 = 5), leaving two vases and the wire. The other player could take two vases and the wire. If they value vases more, that's better for them. My score is lower here, but maybe the other player's proposal would be more complementary. But since my goal is to maximize my own score, taking more vases might be better.} \\
\\ 
\texttt{Wait, but the key is that the proposals need to be complementary. If I take the chair and two vases, the other player can only get one vase and the wire. If their proposal requires the chair, which I have, then our proposals would conflict. But since I don't know their values, maybe it's safer to take the chair and one vase, leaving two vases and the wire. That way, even if the other player wants the chair, they can't get it, but maybe they prefer the wire and vases. Alternatively, if they value the wire highly, taking the wire might be better for them.} \\
\\ 
\texttt{Another angle: The wire is only one, so} \\
            }
        }
    }
    & & \\ \\

    \theutterance \stepcounter{utterance}  
    & & & \multicolumn{2}{p{0.3\linewidth}}{
        \cellcolor[rgb]{0.9,0.9,0.9}{
            \makecell[{{p{\linewidth}}}]{
                \texttt{\tiny{[GM$|$GM]}}
                \texttt{continue} \\
            }
        }
    }
    & & \\ \\

    \theutterance \stepcounter{utterance}  
    & & & \multicolumn{2}{p{0.3\linewidth}}{
        \cellcolor[rgb]{0.9,0.9,0.9}{
            \makecell[{{p{\linewidth}}}]{
                \texttt{\tiny{[GM$|$GM]}}
                \texttt{<think>} \\
\texttt{Okay, let's try to figure out the best way to approach this. The user has to divide 3 vases, 1 chair, and 1 wire. Their values are vase: 2, chair: 3, wire: 1. The other player's values are secret, so I need to make a proposal that's likely to be complementary regardless of what the other player wants.} \\
\\ 
\texttt{First, the chair is the highest value item for the user (3 points). If I can get the chair, that's a big boost. But the chair is unique, so if I take it, the other player can't. The wire is only worth 1 to me, so maybe I should leave that for them if they value it more. The vases are three, so there's more flexibility there.} \\
\\ 
\texttt{If I take the chair and two vases, that's 3 + 2*2 = 7 points. Then the other player would have one vase and the wire. If their values for vase and wire are such that they can get a positive score, that's good. But if they value the chair more, maybe they'd want it, but since I took it, their proposal might conflict. However, since I don't know their values, I have to assume they might want the wire or some vases.} \\
\\ 
\texttt{Alternatively, if I take the chair and one vase (3 + 2 = 5), leaving two vases and the wire. The other player could take two vases and the wire. If they value vases more, that's better for them. My score is lower here, but maybe the other player's proposal would be more complementary. But since my goal is to maximize my own score, taking more vases might be better.} \\
\\ 
\texttt{Wait, but the key is that the proposals need to be complementary. If I take the chair and two vases, the other player can only get one vase and the wire. If their proposal requires the chair, which I have, then our proposals would conflict. But since I don't know their values, maybe it's safer to take the chair and one vase, leaving two vases and the wire. That way, even if the other player wants the chair, they can't get it, but maybe they prefer the wire and vases. Alternatively, if they value the wire highly, taking the wire might be better for them.} \\
\\ 
\texttt{Another angle: The wire is only one, so} \\
            }
        }
    }
    & & \\ \\

    \theutterance \stepcounter{utterance}  
    & & \multicolumn{4}{p{0.6\linewidth}}{
        \cellcolor[rgb]{0.9,0.9,0.9}{
            \makecell[{{p{\linewidth}}}]{
                \texttt{\tiny{[GM$\rangle$P2]}}
                \texttt{<think>} \\
\texttt{Okay, let's try to figure out the best way to approach this. The user has to divide 3 vases, 1 chair, and 1 wire. Their values are vase: 2, chair: 3, wire: 1. The other player's values are secret, so I need to make a proposal that's likely to be complementary regardless of what the other player wants.} \\
\\ 
\texttt{First, the chair is the highest value item for the user (3 points). If I can get the chair, that's a big boost. But the chair is unique, so if I take it, the other player can't. The wire is only worth 1 to me, so maybe I should leave that for them if they value it more. The vases are three, so there's more flexibility there.} \\
\\ 
\texttt{If I take the chair and two vases, that's 3 + 2*2 = 7 points. Then the other player would have one vase and the wire. If their values for vase and wire are such that they can get a positive score, that's good. But if they value the chair more, maybe they'd want it, but since I took it, their proposal might conflict. However, since I don't know their values, I have to assume they might want the wire or some vases.} \\
\\ 
\texttt{Alternatively, if I take the chair and one vase (3 + 2 = 5), leaving two vases and the wire. The other player could take two vases and the wire. If they value vases more, that's better for them. My score is lower here, but maybe the other player's proposal would be more complementary. But since my goal is to maximize my own score, taking more vases might be better.} \\
\\ 
\texttt{Wait, but the key is that the proposals need to be complementary. If I take the chair and two vases, the other player can only get one vase and the wire. If their proposal requires the chair, which I have, then our proposals would conflict. But since I don't know their values, maybe it's safer to take the chair and one vase, leaving two vases and the wire. That way, even if the other player wants the chair, they can't get it, but maybe they prefer the wire and vases. Alternatively, if they value the wire highly, taking the wire might be better for them.} \\
\\ 
\texttt{Another angle: The wire is only one, so} \\
            }
        }
    }
    & \\ \\

    \theutterance \stepcounter{utterance}  
    & & & \multicolumn{4}{p{0.6\linewidth}}{
        \cellcolor[rgb]{0.9,0.9,0.9}{
            \makecell[{{p{\linewidth}}}]{
                \texttt{\tiny{[GM$\langle$P2]}}
                \texttt{<think>} \\
\texttt{Okay, let's try to figure out the best way to approach this. The user has to divide 3 vases, 1 chair, and 1 wire. Their values are vase: 2, chair: 3, wire: 1. The other player's values are secret, so I need to make a proposal that's likely to be complementary regardless of what the other player wants.} \\
\\ 
\texttt{First, the chair is the highest value item for the user (3 points). If I can get the chair, that's a big boost. But the chair is unique, so if I take it, the other player can't. The wire is only worth 1 to me, so maybe I should leave that for them if they value it more. The vases are three, so there's more flexibility there.} \\
\\ 
\texttt{If I take the chair and two vases, that's 3 + 2*2 = 7 points. Then the other player would have one vase and the wire. If their values for vase and wire are such that they can get a positive score, that's good. But if they value the chair more, maybe they'd want it, but since I took it, their proposal might conflict. However, since I don't know their values, I have to assume they might want the wire or some vases.} \\
\\ 
\texttt{Alternatively, if I take the chair and one vase (3 + 2 = 5), leaving two vases and the wire. The other player could take two vases and the wire. If they value vases more, that's better for them. My score is lower here, but maybe the other player's proposal would be more complementary. But since my goal is to maximize my own score, taking more vases might be better.} \\
\\ 
\texttt{Wait, but the key is that the proposals need to be complementary. If I take the chair and two vases, the other player can only get one vase and the wire. If their proposal requires the chair, which I have, then our proposals would conflict. But since I don't know their values, maybe it's safer to take the chair and one vase, leaving two vases and the wire. That way, even if the other player wants the chair, they can't get it, but maybe they prefer the wire and vases. Alternatively, if they value the wire highly, taking the wire might be better for them.} \\
\\ 
\texttt{Another angle: The wire is only one, so} \\
            }
        }
    }
     \\ \\

    \theutterance \stepcounter{utterance}  
    & & & \multicolumn{2}{p{0.3\linewidth}}{
        \cellcolor[rgb]{0.9,0.9,0.9}{
            \makecell[{{p{\linewidth}}}]{
                \texttt{\tiny{[GM$|$GM]}}
                \texttt{continue} \\
            }
        }
    }
    & & \\ \\

    \theutterance \stepcounter{utterance}  
    & & & \multicolumn{2}{p{0.3\linewidth}}{
        \cellcolor[rgb]{0.9,0.9,0.9}{
            \makecell[{{p{\linewidth}}}]{
                \texttt{\tiny{[GM$|$GM]}}
                \texttt{<think>} \\
\texttt{Okay, let's try to figure out the best way to approach this. The user has to divide 3 vases, 1 chair, and 1 wire. Their values are vase: 2, chair: 3, wire: 1. The other player's values are secret, so I need to make a proposal that's likely to be complementary regardless of what the other player wants.} \\
\\ 
\texttt{First, the chair is the highest value item for the user (3 points). If I can get the chair, that's a big boost. But the chair is unique, so if I take it, the other player can't. The wire is only worth 1 to me, so maybe I should leave that for them if they value it more. The vases are three, so there's more flexibility there.} \\
\\ 
\texttt{If I take the chair and two vases, that's 3 + 2*2 = 7 points. Then the other player would have one vase and the wire. If their values for vase and wire are such that they can get a positive score, that's good. But if they value the chair more, maybe they'd want it, but since I took it, their proposal might conflict. However, since I don't know their values, I have to assume they might want the wire or some vases.} \\
\\ 
\texttt{Alternatively, if I take the chair and one vase (3 + 2 = 5), leaving two vases and the wire. The other player could take two vases and the wire. If they value vases more, that's better for them. My score is lower here, but maybe the other player's proposal would be more complementary. But since my goal is to maximize my own score, taking more vases might be better.} \\
\\ 
\texttt{Wait, but the key is that the proposals need to be complementary. If I take the chair and two vases, the other player can only get one vase and the wire. If their proposal requires the chair, which I have, then our proposals would conflict. But since I don't know their values, maybe it's safer to take the chair and one vase, leaving two vases and the wire. That way, even if the other player wants the chair, they can't get it, but maybe they prefer the wire and vases. Alternatively, if they value the wire highly, taking the wire might be better for them.} \\
\\ 
\texttt{Another angle: The wire is only one, so} \\
            }
        }
    }
    & & \\ \\

    \theutterance \stepcounter{utterance}  
    & & \multicolumn{4}{p{0.6\linewidth}}{
        \cellcolor[rgb]{0.9,0.9,0.9}{
            \makecell[{{p{\linewidth}}}]{
                \texttt{\tiny{[P1$\langle$GM]}}
                \texttt{<think>} \\
\texttt{Okay, let's try to figure out the best way to approach this. The user has to divide 3 vases, 1 chair, and 1 wire. Their values are vase: 2, chair: 3, wire: 1. The other player's values are secret, so I need to make a proposal that's likely to be complementary regardless of what the other player wants.} \\
\\ 
\texttt{First, the chair is the highest value item for the user (3 points). If I can get the chair, that's a big boost. But the chair is unique, so if I take it, the other player can't. The wire is only worth 1 to me, so maybe I should leave that for them if they value it more. The vases are three, so there's more flexibility there.} \\
\\ 
\texttt{If I take the chair and two vases, that's 3 + 2*2 = 7 points. Then the other player would have one vase and the wire. If their values for vase and wire are such that they can get a positive score, that's good. But if they value the chair more, maybe they'd want it, but since I took it, their proposal might conflict. However, since I don't know their values, I have to assume they might want the wire or some vases.} \\
\\ 
\texttt{Alternatively, if I take the chair and one vase (3 + 2 = 5), leaving two vases and the wire. The other player could take two vases and the wire. If they value vases more, that's better for them. My score is lower here, but maybe the other player's proposal would be more complementary. But since my goal is to maximize my own score, taking more vases might be better.} \\
\\ 
\texttt{Wait, but the key is that the proposals need to be complementary. If I take the chair and two vases, the other player can only get one vase and the wire. If their proposal requires the chair, which I have, then our proposals would conflict. But since I don't know their values, maybe it's safer to take the chair and one vase, leaving two vases and the wire. That way, even if the other player wants the chair, they can't get it, but maybe they prefer the wire and vases. Alternatively, if they value the wire highly, taking the wire might be better for them.} \\
\\ 
\texttt{Another angle: The wire is only one, so} \\
            }
        }
    }
    & \\ \\

    \theutterance \stepcounter{utterance}  
    & \multicolumn{4}{p{0.6\linewidth}}{
        \cellcolor[rgb]{0.9,0.9,0.9}{
            \makecell[{{p{\linewidth}}}]{
                \texttt{\tiny{[P1$\rangle$GM]}}
                \texttt{<think>} \\
\texttt{Okay, let's try to figure out the best way to approach this. The user has to divide 3 vases, 1 chair, and 1 wire. Their values are vase: 2, chair: 3, wire: 1. The other player's values are secret, so I need to make a proposal that's likely to be complementary regardless of what the other player wants.} \\
\\ 
\texttt{First, the chair is the highest value item for the user (3 points). If I can get the chair, that's a big boost. But the chair is unique, so if I take it, the other player can't. The wire is only worth 1 to me, so maybe I should leave that for them if they value it more. The vases are three, so there's more flexibility there.} \\
\\ 
\texttt{If I take the chair and two vases, that's 3 + 2*2 = 7 points. Then the other player would have one vase and the wire. If their values for vase and wire are such that they can get a positive score, that's good. But if they value the chair more, maybe they'd want it, but since I took it, their proposal might conflict. However, since I don't know their values, I have to assume they might want the wire or some vases.} \\
\\ 
\texttt{Alternatively, if I take the chair and one vase (3 + 2 = 5), leaving two vases and the wire. The other player could take two vases and the wire. If they value vases more, that's better for them. My score is lower here, but maybe the other player's proposal would be more complementary. But since my goal is to maximize my own score, taking more vases might be better.} \\
\\ 
\texttt{Wait, but the key is that the proposals need to be complementary. If I take the chair and two vases, the other player can only get one vase and the wire. If their proposal requires the chair, which I have, then our proposals would conflict. But since I don't know their values, maybe it's safer to take the chair and one vase, leaving two vases and the wire. That way, even if the other player wants the chair, they can't get it, but maybe they prefer the wire and vases. Alternatively, if they value the wire highly, taking the wire might be better for them.} \\
\\ 
\texttt{Another angle: The wire is only one, so} \\
            }
        }
    }
    & & \\ \\

    \theutterance \stepcounter{utterance}  
    & & & \multicolumn{2}{p{0.3\linewidth}}{
        \cellcolor[rgb]{0.9,0.9,0.9}{
            \makecell[{{p{\linewidth}}}]{
                \texttt{\tiny{[GM$|$GM]}}
                \texttt{continue} \\
            }
        }
    }
    & & \\ \\

    \theutterance \stepcounter{utterance}  
    & & & \multicolumn{2}{p{0.3\linewidth}}{
        \cellcolor[rgb]{0.9,0.9,0.9}{
            \makecell[{{p{\linewidth}}}]{
                \texttt{\tiny{[GM$|$GM]}}
                \texttt{<think>} \\
\texttt{Okay, let's try to figure out the best way to approach this. The user has to divide 3 vases, 1 chair, and 1 wire. Their values are vase: 2, chair: 3, wire: 1. The other player's values are secret, so I need to make a proposal that's likely to be complementary regardless of what the other player wants.} \\
\\ 
\texttt{First, the chair is the highest value item for the user (3 points). If I can get the chair, that's a big boost. But the chair is unique, so if I take it, the other player can't. The wire is only worth 1 to me, so maybe I should leave that for them if they value it more. The vases are three, so there's more flexibility there.} \\
\\ 
\texttt{If I take the chair and two vases, that's 3 + 2*2 = 7 points. Then the other player would have one vase and the wire. If their values for vase and wire are such that they can get a positive score, that's good. But if they value the chair more, maybe they'd want it, but since I took it, their proposal might conflict. However, since I don't know their values, I have to assume they might want the wire or some vases.} \\
\\ 
\texttt{Alternatively, if I take the chair and one vase (3 + 2 = 5), leaving two vases and the wire. The other player could take two vases and the wire. If they value vases more, that's better for them. My score is lower here, but maybe the other player's proposal would be more complementary. But since my goal is to maximize my own score, taking more vases might be better.} \\
\\ 
\texttt{Wait, but the key is that the proposals need to be complementary. If I take the chair and two vases, the other player can only get one vase and the wire. If their proposal requires the chair, which I have, then our proposals would conflict. But since I don't know their values, maybe it's safer to take the chair and one vase, leaving two vases and the wire. That way, even if the other player wants the chair, they can't get it, but maybe they prefer the wire and vases. Alternatively, if they value the wire highly, taking the wire might be better for them.} \\
\\ 
\texttt{Another angle: The wire is only one, so} \\
            }
        }
    }
    & & \\ \\

    \theutterance \stepcounter{utterance}  
    & & \multicolumn{4}{p{0.6\linewidth}}{
        \cellcolor[rgb]{0.9,0.9,0.9}{
            \makecell[{{p{\linewidth}}}]{
                \texttt{\tiny{[GM$\rangle$P2]}}
                \texttt{<think>} \\
\texttt{Okay, let's try to figure out the best way to approach this. The user has to divide 3 vases, 1 chair, and 1 wire. Their values are vase: 2, chair: 3, wire: 1. The other player's values are secret, so I need to make a proposal that's likely to be complementary regardless of what the other player wants.} \\
\\ 
\texttt{First, the chair is the highest value item for the user (3 points). If I can get the chair, that's a big boost. But the chair is unique, so if I take it, the other player can't. The wire is only worth 1 to me, so maybe I should leave that for them if they value it more. The vases are three, so there's more flexibility there.} \\
\\ 
\texttt{If I take the chair and two vases, that's 3 + 2*2 = 7 points. Then the other player would have one vase and the wire. If their values for vase and wire are such that they can get a positive score, that's good. But if they value the chair more, maybe they'd want it, but since I took it, their proposal might conflict. However, since I don't know their values, I have to assume they might want the wire or some vases.} \\
\\ 
\texttt{Alternatively, if I take the chair and one vase (3 + 2 = 5), leaving two vases and the wire. The other player could take two vases and the wire. If they value vases more, that's better for them. My score is lower here, but maybe the other player's proposal would be more complementary. But since my goal is to maximize my own score, taking more vases might be better.} \\
\\ 
\texttt{Wait, but the key is that the proposals need to be complementary. If I take the chair and two vases, the other player can only get one vase and the wire. If their proposal requires the chair, which I have, then our proposals would conflict. But since I don't know their values, maybe it's safer to take the chair and one vase, leaving two vases and the wire. That way, even if the other player wants the chair, they can't get it, but maybe they prefer the wire and vases. Alternatively, if they value the wire highly, taking the wire might be better for them.} \\
\\ 
\texttt{Another angle: The wire is only one, so} \\
            }
        }
    }
    & \\ \\

    \theutterance \stepcounter{utterance}  
    & & & \multicolumn{4}{p{0.6\linewidth}}{
        \cellcolor[rgb]{0.9,0.9,0.9}{
            \makecell[{{p{\linewidth}}}]{
                \texttt{\tiny{[GM$\langle$P2]}}
                \texttt{<think>} \\
\texttt{Okay, let's try to figure out the best way to approach this. The user has to divide 3 vases, 1 chair, and 1 wire. Their values are vase: 2, chair: 3, wire: 1. The other player's values are secret, so I need to make a proposal that's likely to be complementary regardless of what the other player wants.} \\
\\ 
\texttt{First, the chair is the highest value item for the user (3 points). If I can get the chair, that's a big boost. But the chair is unique, so if I take it, the other player can't. The wire is only worth 1 to me, so maybe I should leave that for them if they value it more. The vases are three, so there's more flexibility there.} \\
\\ 
\texttt{If I take the chair and two vases, that's 3 + 2*2 = 7 points. Then the other player would have one vase and the wire. If their values for vase and wire are such that they can get a positive score, that's good. But if they value the chair more, maybe they'd want it, but since I took it, their proposal might conflict. However, since I don't know their values, I have to assume they might want the wire or some vases.} \\
\\ 
\texttt{Alternatively, if I take the chair and one vase (3 + 2 = 5), leaving two vases and the wire. The other player could take two vases and the wire. If they value vases more, that's better for them. My score is lower here, but maybe the other player's proposal would be more complementary. But since my goal is to maximize my own score, taking more vases might be better.} \\
\\ 
\texttt{Wait, but the key is that the proposals need to be complementary. If I take the chair and two vases, the other player can only get one vase and the wire. If their proposal requires the chair, which I have, then our proposals would conflict. But since I don't know their values, maybe it's safer to take the chair and one vase, leaving two vases and the wire. That way, even if the other player wants the chair, they can't get it, but maybe they prefer the wire and vases. Alternatively, if they value the wire highly, taking the wire might be better for them.} \\
\\ 
\texttt{Another angle: The wire is only one, so} \\
            }
        }
    }
     \\ \\

    \theutterance \stepcounter{utterance}  
    & & & \multicolumn{2}{p{0.3\linewidth}}{
        \cellcolor[rgb]{0.9,0.9,0.9}{
            \makecell[{{p{\linewidth}}}]{
                \texttt{\tiny{[GM$|$GM]}}
                \texttt{continue} \\
            }
        }
    }
    & & \\ \\

    \theutterance \stepcounter{utterance}  
    & & & \multicolumn{2}{p{0.3\linewidth}}{
        \cellcolor[rgb]{0.9,0.9,0.9}{
            \makecell[{{p{\linewidth}}}]{
                \texttt{\tiny{[GM$|$GM]}}
                \texttt{<think>} \\
\texttt{Okay, let's try to figure out the best way to approach this. The user has to divide 3 vases, 1 chair, and 1 wire. Their values are vase: 2, chair: 3, wire: 1. The other player's values are secret, so I need to make a proposal that's likely to be complementary regardless of what the other player wants.} \\
\\ 
\texttt{First, the chair is the highest value item for the user (3 points). If I can get the chair, that's a big boost. But the chair is unique, so if I take it, the other player can't. The wire is only worth 1 to me, so maybe I should leave that for them if they value it more. The vases are three, so there's more flexibility there.} \\
\\ 
\texttt{If I take the chair and two vases, that's 3 + 2*2 = 7 points. Then the other player would have one vase and the wire. If their values for vase and wire are such that they can get a positive score, that's good. But if they value the chair more, maybe they'd want it, but since I took it, their proposal might conflict. However, since I don't know their values, I have to assume they might want the wire or some vases.} \\
\\ 
\texttt{Alternatively, if I take the chair and one vase (3 + 2 = 5), leaving two vases and the wire. The other player could take two vases and the wire. If they value vases more, that's better for them. My score is lower here, but maybe the other player's proposal would be more complementary. But since my goal is to maximize my own score, taking more vases might be better.} \\
\\ 
\texttt{Wait, but the key is that the proposals need to be complementary. If I take the chair and two vases, the other player can only get one vase and the wire. If their proposal requires the chair, which I have, then our proposals would conflict. But since I don't know their values, maybe it's safer to take the chair and one vase, leaving two vases and the wire. That way, even if the other player wants the chair, they can't get it, but maybe they prefer the wire and vases. Alternatively, if they value the wire highly, taking the wire might be better for them.} \\
\\ 
\texttt{Another angle: The wire is only one, so} \\
            }
        }
    }
    & & \\ \\

    \theutterance \stepcounter{utterance}  
    & & \multicolumn{4}{p{0.6\linewidth}}{
        \cellcolor[rgb]{0.9,0.9,0.9}{
            \makecell[{{p{\linewidth}}}]{
                \texttt{\tiny{[P1$\langle$GM]}}
                \texttt{<think>} \\
\texttt{Okay, let's try to figure out the best way to approach this. The user has to divide 3 vases, 1 chair, and 1 wire. Their values are vase: 2, chair: 3, wire: 1. The other player's values are secret, so I need to make a proposal that's likely to be complementary regardless of what the other player wants.} \\
\\ 
\texttt{First, the chair is the highest value item for the user (3 points). If I can get the chair, that's a big boost. But the chair is unique, so if I take it, the other player can't. The wire is only worth 1 to me, so maybe I should leave that for them if they value it more. The vases are three, so there's more flexibility there.} \\
\\ 
\texttt{If I take the chair and two vases, that's 3 + 2*2 = 7 points. Then the other player would have one vase and the wire. If their values for vase and wire are such that they can get a positive score, that's good. But if they value the chair more, maybe they'd want it, but since I took it, their proposal might conflict. However, since I don't know their values, I have to assume they might want the wire or some vases.} \\
\\ 
\texttt{Alternatively, if I take the chair and one vase (3 + 2 = 5), leaving two vases and the wire. The other player could take two vases and the wire. If they value vases more, that's better for them. My score is lower here, but maybe the other player's proposal would be more complementary. But since my goal is to maximize my own score, taking more vases might be better.} \\
\\ 
\texttt{Wait, but the key is that the proposals need to be complementary. If I take the chair and two vases, the other player can only get one vase and the wire. If their proposal requires the chair, which I have, then our proposals would conflict. But since I don't know their values, maybe it's safer to take the chair and one vase, leaving two vases and the wire. That way, even if the other player wants the chair, they can't get it, but maybe they prefer the wire and vases. Alternatively, if they value the wire highly, taking the wire might be better for them.} \\
\\ 
\texttt{Another angle: The wire is only one, so} \\
            }
        }
    }
    & \\ \\

    \theutterance \stepcounter{utterance}  
    & \multicolumn{4}{p{0.6\linewidth}}{
        \cellcolor[rgb]{0.9,0.9,0.9}{
            \makecell[{{p{\linewidth}}}]{
                \texttt{\tiny{[P1$\rangle$GM]}}
                \texttt{<think>} \\
\texttt{Okay, let's try to figure out the best way to approach this. The user has to divide 3 vases, 1 chair, and 1 wire. Their values are vase: 2, chair: 3, wire: 1. The other player's values are secret, so I need to make a proposal that's likely to be complementary regardless of what the other player wants.} \\
\\ 
\texttt{First, the chair is the highest value item for the user (3 points). If I can get the chair, that's a big boost. But the chair is unique, so if I take it, the other player can't. The wire is only worth 1 to me, so maybe I should leave that for them if they value it more. The vases are three, so there's more flexibility there.} \\
\\ 
\texttt{If I take the chair and two vases, that's 3 + 2*2 = 7 points. Then the other player would have one vase and the wire. If their values for vase and wire are such that they can get a positive score, that's good. But if they value the chair more, maybe they'd want it, but since I took it, their proposal might conflict. However, since I don't know their values, I have to assume they might want the wire or some vases.} \\
\\ 
\texttt{Alternatively, if I take the chair and one vase (3 + 2 = 5), leaving two vases and the wire. The other player could take two vases and the wire. If they value vases more, that's better for them. My score is lower here, but maybe the other player's proposal would be more complementary. But since my goal is to maximize my own score, taking more vases might be better.} \\
\\ 
\texttt{Wait, but the key is that the proposals need to be complementary. If I take the chair and two vases, the other player can only get one vase and the wire. If their proposal requires the chair, which I have, then our proposals would conflict. But since I don't know their values, maybe it's safer to take the chair and one vase, leaving two vases and the wire. That way, even if the other player wants the chair, they can't get it, but maybe they prefer the wire and vases. Alternatively, if they value the wire highly, taking the wire might be better for them.} \\
\\ 
\texttt{Another angle: The wire is only one, so} \\
            }
        }
    }
    & & \\ \\

    \theutterance \stepcounter{utterance}  
    & & & \multicolumn{2}{p{0.3\linewidth}}{
        \cellcolor[rgb]{0.9,0.9,0.9}{
            \makecell[{{p{\linewidth}}}]{
                \texttt{\tiny{[GM$|$GM]}}
                \texttt{continue} \\
            }
        }
    }
    & & \\ \\

    \theutterance \stepcounter{utterance}  
    & & & \multicolumn{2}{p{0.3\linewidth}}{
        \cellcolor[rgb]{0.9,0.9,0.9}{
            \makecell[{{p{\linewidth}}}]{
                \texttt{\tiny{[GM$|$GM]}}
                \texttt{<think>} \\
\texttt{Okay, let's try to figure out the best way to approach this. The user has to divide 3 vases, 1 chair, and 1 wire. Their values are vase: 2, chair: 3, wire: 1. The other player's values are secret, so I need to make a proposal that's likely to be complementary regardless of what the other player wants.} \\
\\ 
\texttt{First, the chair is the highest value item for the user (3 points). If I can get the chair, that's a big boost. But the chair is unique, so if I take it, the other player can't. The wire is only worth 1 to me, so maybe I should leave that for them if they value it more. The vases are three, so there's more flexibility there.} \\
\\ 
\texttt{If I take the chair and two vases, that's 3 + 2*2 = 7 points. Then the other player would have one vase and the wire. If their values for vase and wire are such that they can get a positive score, that's good. But if they value the chair more, maybe they'd want it, but since I took it, their proposal might conflict. However, since I don't know their values, I have to assume they might want the wire or some vases.} \\
\\ 
\texttt{Alternatively, if I take the chair and one vase (3 + 2 = 5), leaving two vases and the wire. The other player could take two vases and the wire. If they value vases more, that's better for them. My score is lower here, but maybe the other player's proposal would be more complementary. But since my goal is to maximize my own score, taking more vases might be better.} \\
\\ 
\texttt{Wait, but the key is that the proposals need to be complementary. If I take the chair and two vases, the other player can only get one vase and the wire. If their proposal requires the chair, which I have, then our proposals would conflict. But since I don't know their values, maybe it's safer to take the chair and one vase, leaving two vases and the wire. That way, even if the other player wants the chair, they can't get it, but maybe they prefer the wire and vases. Alternatively, if they value the wire highly, taking the wire might be better for them.} \\
\\ 
\texttt{Another angle: The wire is only one, so} \\
            }
        }
    }
    & & \\ \\

    \theutterance \stepcounter{utterance}  
    & & \multicolumn{4}{p{0.6\linewidth}}{
        \cellcolor[rgb]{0.9,0.9,0.9}{
            \makecell[{{p{\linewidth}}}]{
                \texttt{\tiny{[GM$\rangle$P2]}}
                \texttt{<think>} \\
\texttt{Okay, let's try to figure out the best way to approach this. The user has to divide 3 vases, 1 chair, and 1 wire. Their values are vase: 2, chair: 3, wire: 1. The other player's values are secret, so I need to make a proposal that's likely to be complementary regardless of what the other player wants.} \\
\\ 
\texttt{First, the chair is the highest value item for the user (3 points). If I can get the chair, that's a big boost. But the chair is unique, so if I take it, the other player can't. The wire is only worth 1 to me, so maybe I should leave that for them if they value it more. The vases are three, so there's more flexibility there.} \\
\\ 
\texttt{If I take the chair and two vases, that's 3 + 2*2 = 7 points. Then the other player would have one vase and the wire. If their values for vase and wire are such that they can get a positive score, that's good. But if they value the chair more, maybe they'd want it, but since I took it, their proposal might conflict. However, since I don't know their values, I have to assume they might want the wire or some vases.} \\
\\ 
\texttt{Alternatively, if I take the chair and one vase (3 + 2 = 5), leaving two vases and the wire. The other player could take two vases and the wire. If they value vases more, that's better for them. My score is lower here, but maybe the other player's proposal would be more complementary. But since my goal is to maximize my own score, taking more vases might be better.} \\
\\ 
\texttt{Wait, but the key is that the proposals need to be complementary. If I take the chair and two vases, the other player can only get one vase and the wire. If their proposal requires the chair, which I have, then our proposals would conflict. But since I don't know their values, maybe it's safer to take the chair and one vase, leaving two vases and the wire. That way, even if the other player wants the chair, they can't get it, but maybe they prefer the wire and vases. Alternatively, if they value the wire highly, taking the wire might be better for them.} \\
\\ 
\texttt{Another angle: The wire is only one, so} \\
            }
        }
    }
    & \\ \\

    \theutterance \stepcounter{utterance}  
    & & & \multicolumn{4}{p{0.6\linewidth}}{
        \cellcolor[rgb]{0.9,0.9,0.9}{
            \makecell[{{p{\linewidth}}}]{
                \texttt{\tiny{[GM$\langle$P2]}}
                \texttt{<think>} \\
\texttt{Okay, let's try to figure out the best way to approach this. The user has to divide 3 vases, 1 chair, and 1 wire. Their values are vase: 2, chair: 3, wire: 1. The other player's values are secret, so I need to make a proposal that's likely to be complementary regardless of what the other player wants.} \\
\\ 
\texttt{First, the chair is the highest value item for the user (3 points). If I can get the chair, that's a big boost. But the chair is unique, so if I take it, the other player can't. The wire is only worth 1 to me, so maybe I should leave that for them if they value it more. The vases are three, so there's more flexibility there.} \\
\\ 
\texttt{If I take the chair and two vases, that's 3 + 2*2 = 7 points. Then the other player would have one vase and the wire. If their values for vase and wire are such that they can get a positive score, that's good. But if they value the chair more, maybe they'd want it, but since I took it, their proposal might conflict. However, since I don't know their values, I have to assume they might want the wire or some vases.} \\
\\ 
\texttt{Alternatively, if I take the chair and one vase (3 + 2 = 5), leaving two vases and the wire. The other player could take two vases and the wire. If they value vases more, that's better for them. My score is lower here, but maybe the other player's proposal would be more complementary. But since my goal is to maximize my own score, taking more vases might be better.} \\
\\ 
\texttt{Wait, but the key is that the proposals need to be complementary. If I take the chair and two vases, the other player can only get one vase and the wire. If their proposal requires the chair, which I have, then our proposals would conflict. But since I don't know their values, maybe it's safer to take the chair and one vase, leaving two vases and the wire. That way, even if the other player wants the chair, they can't get it, but maybe they prefer the wire and vases. Alternatively, if they value the wire highly, taking the wire might be better for them.} \\
\\ 
\texttt{Another angle: The wire is only one, so} \\
            }
        }
    }
     \\ \\

    \theutterance \stepcounter{utterance}  
    & & & \multicolumn{2}{p{0.3\linewidth}}{
        \cellcolor[rgb]{0.9,0.9,0.9}{
            \makecell[{{p{\linewidth}}}]{
                \texttt{\tiny{[GM$|$GM]}}
                \texttt{continue} \\
            }
        }
    }
    & & \\ \\

    \theutterance \stepcounter{utterance}  
    & & & \multicolumn{2}{p{0.3\linewidth}}{
        \cellcolor[rgb]{0.9,0.9,0.9}{
            \makecell[{{p{\linewidth}}}]{
                \texttt{\tiny{[GM$|$GM]}}
                \texttt{<think>} \\
\texttt{Okay, let's try to figure out the best way to approach this. The user has to divide 3 vases, 1 chair, and 1 wire. Their values are vase: 2, chair: 3, wire: 1. The other player's values are secret, so I need to make a proposal that's likely to be complementary regardless of what the other player wants.} \\
\\ 
\texttt{First, the chair is the highest value item for the user (3 points). If I can get the chair, that's a big boost. But the chair is unique, so if I take it, the other player can't. The wire is only worth 1 to me, so maybe I should leave that for them if they value it more. The vases are three, so there's more flexibility there.} \\
\\ 
\texttt{If I take the chair and two vases, that's 3 + 2*2 = 7 points. Then the other player would have one vase and the wire. If their values for vase and wire are such that they can get a positive score, that's good. But if they value the chair more, maybe they'd want it, but since I took it, their proposal might conflict. However, since I don't know their values, I have to assume they might want the wire or some vases.} \\
\\ 
\texttt{Alternatively, if I take the chair and one vase (3 + 2 = 5), leaving two vases and the wire. The other player could take two vases and the wire. If they value vases more, that's better for them. My score is lower here, but maybe the other player's proposal would be more complementary. But since my goal is to maximize my own score, taking more vases might be better.} \\
\\ 
\texttt{Wait, but the key is that the proposals need to be complementary. If I take the chair and two vases, the other player can only get one vase and the wire. If their proposal requires the chair, which I have, then our proposals would conflict. But since I don't know their values, maybe it's safer to take the chair and one vase, leaving two vases and the wire. That way, even if the other player wants the chair, they can't get it, but maybe they prefer the wire and vases. Alternatively, if they value the wire highly, taking the wire might be better for them.} \\
\\ 
\texttt{Another angle: The wire is only one, so} \\
            }
        }
    }
    & & \\ \\

    \theutterance \stepcounter{utterance}  
    & & \multicolumn{4}{p{0.6\linewidth}}{
        \cellcolor[rgb]{0.9,0.9,0.9}{
            \makecell[{{p{\linewidth}}}]{
                \texttt{\tiny{[P1$\langle$GM]}}
                \texttt{<think>} \\
\texttt{Okay, let's try to figure out the best way to approach this. The user has to divide 3 vases, 1 chair, and 1 wire. Their values are vase: 2, chair: 3, wire: 1. The other player's values are secret, so I need to make a proposal that's likely to be complementary regardless of what the other player wants.} \\
\\ 
\texttt{First, the chair is the highest value item for the user (3 points). If I can get the chair, that's a big boost. But the chair is unique, so if I take it, the other player can't. The wire is only worth 1 to me, so maybe I should leave that for them if they value it more. The vases are three, so there's more flexibility there.} \\
\\ 
\texttt{If I take the chair and two vases, that's 3 + 2*2 = 7 points. Then the other player would have one vase and the wire. If their values for vase and wire are such that they can get a positive score, that's good. But if they value the chair more, maybe they'd want it, but since I took it, their proposal might conflict. However, since I don't know their values, I have to assume they might want the wire or some vases.} \\
\\ 
\texttt{Alternatively, if I take the chair and one vase (3 + 2 = 5), leaving two vases and the wire. The other player could take two vases and the wire. If they value vases more, that's better for them. My score is lower here, but maybe the other player's proposal would be more complementary. But since my goal is to maximize my own score, taking more vases might be better.} \\
\\ 
\texttt{Wait, but the key is that the proposals need to be complementary. If I take the chair and two vases, the other player can only get one vase and the wire. If their proposal requires the chair, which I have, then our proposals would conflict. But since I don't know their values, maybe it's safer to take the chair and one vase, leaving two vases and the wire. That way, even if the other player wants the chair, they can't get it, but maybe they prefer the wire and vases. Alternatively, if they value the wire highly, taking the wire might be better for them.} \\
\\ 
\texttt{Another angle: The wire is only one, so} \\
\\ 
\\ 
\texttt{The time is up. It is now your turn to submit a secret proposal.} \\
            }
        }
    }
    & \\ \\

    \theutterance \stepcounter{utterance}  
    & \multicolumn{4}{p{0.6\linewidth}}{
        \cellcolor[rgb]{0.9,0.9,0.9}{
            \makecell[{{p{\linewidth}}}]{
                \texttt{\tiny{[P1$\rangle$GM]}}
                \texttt{<think>} \\
\texttt{Okay, let's try to figure out the best way to approach this. The user has to divide 3 vases, 1 chair, and 1 wire. Their values are vase: 2, chair: 3, wire: 1. The other player's values are secret, so I need to make a proposal that's likely to be complementary regardless of what the other player wants.} \\
\\ 
\texttt{First, the chair is the highest value item for the user (3 points). If I can get the chair, that's a big boost. But the chair is unique, so if I take it, the other player can't. The wire is only worth 1 to me, so maybe I should leave that for them if they value it more. The vases are three, so there's more flexibility there.} \\
\\ 
\texttt{If I take the chair and two vases, that's 3 + 2*2 = 7 points. Then the other player would have one vase and the wire. If their values for vase and wire are such that they can get a positive score, that's good. But if they value the chair more, maybe they'd want it, but since I took it, their proposal might conflict. However, since I don't know their values, I have to assume they might want the wire or some vases.} \\
\\ 
\texttt{Alternatively, if I take the chair and one vase (3 + 2 = 5), leaving two vases and the wire. The other player could take two vases and the wire. If they value vases more, that's better for them. My score is lower here, but maybe the other player's proposal would be more complementary. But since my goal is to maximize my own score, taking more vases might be better.} \\
\\ 
\texttt{Wait, but the key is that the proposals need to be complementary. If I take the chair and two vases, the other player can only get one vase and the wire. If their proposal requires the chair, which I have, then our proposals would conflict. But since I don't know their values, maybe it's safer to take the chair and one vase, leaving two vases and the wire. That way, even if the other player wants the chair, they can't get it, but maybe they prefer the wire and vases. Alternatively, if they value the wire highly, taking the wire might be better for them.} \\
\\ 
\texttt{Another angle: The wire is only one, so} \\
            }
        }
    }
    & & \\ \\

    \theutterance \stepcounter{utterance}  
    & & & \multicolumn{2}{p{0.3\linewidth}}{
        \cellcolor[rgb]{0.9,0.9,0.9}{
            \makecell[{{p{\linewidth}}}]{
                \texttt{\tiny{[GM$|$GM]}}
                \texttt{continue} \\
            }
        }
    }
    & & \\ \\

    \theutterance \stepcounter{utterance}  
    & & & \multicolumn{2}{p{0.3\linewidth}}{
        \cellcolor[rgb]{0.9,0.9,0.9}{
            \makecell[{{p{\linewidth}}}]{
                \texttt{\tiny{[GM$|$GM]}}
                \texttt{player was instructed to make a proposal} \\
            }
        }
    }
    & & \\ \\

\end{supertabular}
}

\end{document}
